% from ega4-2-fr.tex
\oldpage[IV-2]{204}

Cela étant, soit $f={}^a\varphi:\operatorname{Spec}(B)\to
\operatorname{Spec}(A)$, et supposons que la propriété
$\boldsymbol P(\operatorname{Spec}(B\otimes_A k),k)$ soit vraie. Alors, pour
tout $x\in\operatorname{Spec}(A)$ la propriété
$\boldsymbol P(f^{-1}(x),k(x))$ est vraie.
\end{theorem}

\begin{proof}
Autrement dit, la propriété $\boldsymbol P$ pour la fibre du point fermé de
$\operatorname{Spec}(A)$ entraîne cette même propriété pour toutes les
fibres du morphisme $f$ (c'est-à-dire que $f$ est un
$\boldsymbol P$-morphisme (7.3.1)). Nous procéderons en plusieurs étapes :

\emph{I) Réduction à l'étude des anneaux locaux de la fibre générique.} ---
Appliquons le lemme (7.3.16.2) : il suffit de voir que pour toute
$A$-algèbre finie et intègre $A'$, de corps des fractions $K'$, si l'on
pose $B'=B\otimes_A A'$, tous les anneaux locaux de
$\operatorname{Spec}(B'\otimes_{A'}K')$ vérifient la propriété
$\boldsymbol R$. L'anneau $A'$ (resp. $B'$) est semi-local complet ($B'$
étant une $B$-algèbre finie), donc produit d'anneaux locaux complets; comme
$A'$ est supposé intègre, il est local; tout idéal maximal $\mathfrak n'$ de
$B'$ est au-dessus de l'idéal maximal $\mathfrak n$ de $B$, donc au-dessus
de $\mathfrak m$, et par suite son image réciproque dans $A'$ est l'idéal
maximal $\mathfrak m'$ de $A'$. Comme tout anneau local de
$\operatorname{Spec}(B'\otimes_{A'}K')$ est un anneau local de
$\operatorname{Spec}(B')$, donc de l'un des
$\operatorname{Spec}(B'_{\mathfrak n'})$ (au-dessus du point générique de
$A'$), on voit qu'il suffira de prouver que les anneaux locaux de
\[
\operatorname{Spec}(B'_{\mathfrak n'}\otimes_{A'}K')
\]
possèdent la propriété $\boldsymbol R$. Or $A'$ et
$B'_{\mathfrak n'}$ sont des anneaux locaux noethériens complets; le corps
résiduel de $B'_{\mathfrak n'}$ est extension finie de celui de $B$, donc
aussi de $k$, et \emph{a fortiori} du corps résiduel $k'$ de $A'$; cette
remarque et (2.1.4) montrent que $A'$ et $B'_{\mathfrak n'}$ vérifient les
conditions (i) et (ii) de l'énoncé. D'autre part, si $k''$ est une extension
finie de $k'$, $k''$ est aussi extension finie de $k$, et les anneaux locaux
de
\[
\operatorname{Spec}(B'_{\mathfrak n'}\otimes_{A'}k'')
\]
sont aussi des anneaux locaux de $\operatorname{Spec}(B\otimes_A k'')$;
donc l'hypothèse que
$\boldsymbol P(\operatorname{Spec}(B\otimes_A k),k)$ est vraie entraîne que
\[
\boldsymbol P(\operatorname{Spec}(B'_{\mathfrak n'}\otimes_{A'}k'),k')
\]
est vraie.

On est ainsi ramené au cas où $A$ est de plus supposé intègre, de corps des
fractions $K$, et à prouver que les anneaux locaux de
$\operatorname{Spec}(B\otimes_A K)$ possèdent la propriété
$\boldsymbol R$.

\emph{II) Cas où $\dim(A)=1$.} --- Soit $A'$ la clôture intégrale de $A$;
on sait (0, 23.1.6) que $A'$ est une $A$-algèbre finie et un anneau local
complet; si l'on pose $B'=B\otimes_A A'$, on a
$B\otimes_A K=B'\otimes_{A'}K$. Le même raisonnement que dans I) montre
qu'il suffit, pour tout idéal maximal $\mathfrak n'$ de $B'$, de prouver que
les anneaux locaux de
$B'_{\mathfrak n'}\otimes_{A'}K$ vérifient $\boldsymbol R$; en outre, ce
raisonnement montre aussi que $A'$ et $B'_{\mathfrak n'}$ vérifient les
hypothèses (i) et (ii) de l'énoncé et que
$\boldsymbol P(\operatorname{Spec}(B'_{\mathfrak n'}\otimes_{A'}k'),k')$
est vraie ($k'$ étant le corps résiduel de $A'$). En outre, comme
$\dim(A')=1$ (0, 16.1.5) et que $A'$ est intègre et intégralement clos,
c'est un anneau de valuation discrète complet. On peut donc se borner au
cas où $A$ est déjà un anneau de valuation discrète complet; si alors $u$
est une uniformisante de $A$, le fait que $u$ soit $A$-régulier et que $B$
soit un $A$-module plat entraîne que $t=\varphi(u)$ est un élément
$B$-régulier, qui appartient à l'idéal maximal de $B$. L'hypothèse que
$\boldsymbol P(\operatorname{Spec}(B\otimes_A k),k)$ est vraie entraîne en
particulier que $\boldsymbol R(B/tB)$ est vraie; en vertu de
$(\mathrm R_{\mathrm{IV}})$, $\boldsymbol R(B)$ est donc vraie. Autrement
dit, $U_{\boldsymbol R}(\operatorname{Spec}(B))$ contient le point fermé de
$\operatorname{Spec}(B)$. Mais comme
$U_{\boldsymbol R}(\operatorname{Spec}(B))$ est ouvert en vertu de
$(\mathrm R_{\mathrm{III}})$ et que $\operatorname{Spec}(B)$ est le seul
ensemble ouvert de $\operatorname{Spec}(B)$ contenant le point fermé, tous
les anneaux locaux de $\operatorname{Spec}(B)$ possèdent la propriété
$\boldsymbol R$, et en particulier ceux de
$\operatorname{Spec}(B\otimes_A K)$.

\emph{III) Cas général.} --- Le cas $\dim(A)=0$ est trivial, puisqu'alors
$A=k$; on peut donc se borner au cas où $\dim(A)\geqslant1$. On sait
(6.12.7) que $\operatorname{Spec}(A)$ contient un ouvert non vide $V$ dont
tous les points sont réguliers; comme $\dim(A)\geqslant1$, l'intersection
$V'$ de $V$ et du complémentaire du point fermé de
$\operatorname{Spec}(A)$ est un ouvert non vide, contenant donc le point
générique de $\operatorname{Spec}(A)$. Si nous prouvons que
\[
f^{-1}(V')\subset U_{\boldsymbol R}(\operatorname{Spec}(B)),
\]
la proposition sera démontrée \emph{a fortiori}. Autrement dit, il suffit de
voir que l'ensemble $Z$, intersection de $f^{-1}(V')$ et du complémentaire
de $U_{\boldsymbol R}(\operatorname{Spec}(B))$, est vide. Raisonnons par
l'absurde : comme, en vertu de $(\mathrm R_{\mathrm{III}})$,
$U_{\boldsymbol R}(\operatorname{Spec}(B))$ est ouvert dans
$\operatorname{Spec}(B)$, $Z$ est localement fermé dans
$\operatorname{Spec}(B)$; s'il n'était pas vide, il contiendrait un point
$x$ tel que $\dim(\overline{\{x\}})\leqslant1$ (5.1.10); comme
$f(x)\in V'$ est distinct du point fermé de $\operatorname{Spec}(A)$, $x$
est distinct du point fermé de $\operatorname{Spec}(B)$, autrement dit
$\dim(\overline{\{x\}})=1$. Nous allons prouver que cela est impossible,
autrement dit que, si $\dim(\overline{\{x\}})=1$ et $f(x)\in V'$ alors on a
$\boldsymbol R(\mathscr O_x)$. En d'autres termes, il s'agit de voir que,
si $\mathfrak q$ est un idéal premier de $B$ tel que
$\dim(B/\mathfrak q)=1$ et si $\mathfrak p=\varphi^{-1}(\mathfrak q)$ est
distinct de $\mathfrak m$ et tel que $A_{\mathfrak p}$ soit régulier, alors
on a $\boldsymbol R(B_{\mathfrak q})$. Comme
$\mathfrak p\ne\mathfrak m$, $\mathfrak m(B/\mathfrak q)$ n'est pas réduit
à 0, donc est un idéal de définition de l'anneau local intègre
$B/\mathfrak q$ de dimension 1; d'autre part, en vertu de (i), le corps
résiduel de $B/\mathfrak q$ est une extension finie du corps résiduel $k$
de $A/\mathfrak p$, donc $B/\mathfrak q$ est une
$(A/\mathfrak p)$-algèbre quasi-finie (0$_{\mathrm I}$, 7.4.4). Mais
$A/\mathfrak p$ est complet et $B/\mathfrak q$ est séparé pour la topologie
$\mathfrak m$-préadique (qui est identique à sa topologie d'anneau local);
donc (0$_{\mathrm I}$, 7.4.1), $B/\mathfrak q$ est une
$(A/\mathfrak p)$-algèbre finie. En outre, par définition,
l'homomorphisme $A/\mathfrak p\to B/\mathfrak q$ est injectif, donc (0,
16.1.5) on a
\[
\dim(A/\mathfrak p)=\dim(B/\mathfrak q)=1.
\]
On peut alors appliquer aux anneaux $A/\mathfrak p$ et
$B/\mathfrak pB$ le résultat de II), car les corps résiduels de ces anneaux
locaux sont respectivement ceux de $A$ et de $B$, et $B/\mathfrak pB$ est
un $(A/\mathfrak p)$-module plat; en outre, on a
\[
(B/\mathfrak pB)\otimes_{A/\mathfrak p}k=B\otimes_A k=B/\mathfrak mB.
\]
Par suite, les anneaux locaux de
\[
\operatorname{Spec}((B/\mathfrak pB)
\otimes_{A/\mathfrak p}k(\mathfrak p))
\]
vérifient $\boldsymbol R$. Or, $B_{\mathfrak q}/\mathfrak pB_{\mathfrak q}$
est un des anneaux locaux de
$\operatorname{Spec}((B/\mathfrak pB)
\otimes_{A/\mathfrak p}k(\mathfrak p))$. En outre,
$B_{\mathfrak q}$ est un $A_{\mathfrak p}$-module plat, et
$A_{\mathfrak p}$ est régulier. Or on a le lemme suivant.

